% Packages, notation and theorem environments.
% Notation follows Annex B of the leanVM specification (doc/leanvm/preamble/macros.tex).
\usepackage[T1]{fontenc}
\usepackage[utf8]{inputenc}
\usepackage{lmodern}
\usepackage[margin=1in]{geometry}
\usepackage{microtype}
\usepackage{amsmath,amssymb,amsthm,mathtools}
\usepackage{tabularx,booktabs,longtable,array}
\usepackage{enumitem}
\usepackage{xcolor}
\usepackage{tikz}
\usetikzlibrary{arrows.meta}
\usepackage{listings}
\usepackage[colorlinks=true,linkcolor=blue!50!black,citecolor=blue!50!black,urlcolor=blue!50!black]{hyperref}

\emergencystretch=1.5em

% ragged-right table cells: narrow justified columns set badly
\newcolumntype{L}[1]{>{\raggedright\arraybackslash}p{#1}}
\newcolumntype{Y}{>{\raggedright\arraybackslash}X}

% fields, multilinears
\newcommand{\F}{\mathbb{F}}
\newcommand{\Ftwo}{\mathbb{F}_2}
\newcommand{\K}{K}
\newcommand{\E}{E}
\newcommand{\q}{q}
\newcommand{\mle}[1]{\widetilde{#1}}
\newcommand{\eq}{\operatorname{eq}}
\newcommand{\cube}[1]{\{0,1\}^{#1}}
\newcommand{\inner}[2]{\langle #1,\, #2 \rangle}
\newcommand{\fc}{\chi}
\newcommand{\qeq}{\stackrel{?}{=}}

% codes and the PCS, as in Annex B
\newcommand{\Enc}{\mathrm{Enc}}
\newcommand{\RS}{\mathrm{RS}}
\newcommand{\nv}{M}
\newcommand{\nlev}{\Lambda}
\newcommand{\rate}{\rho}
\newcommand{\rexp}{\nu}
\newcommand{\rc}[1]{c^{(#1)}}
\newcommand{\nq}{\omega}
\newcommand{\blen}{n}
\newcommand{\dom}{\mathcal{D}}
\newcommand{\orc}{Z}
\newcommand{\lst}{\mathcal{L}}
\newcommand{\aset}{\mathcal{A}}
\newcommand{\agr}{t}
\newcommand{\dgr}{\mathsf{d}}
\newcommand{\jl}{\mathsf{L}}
\newcommand{\afa}{\mathsf{a}}
\newcommand{\afb}{\mathsf{b}}
\newcommand{\tabx}[1]{v_{X,#1}}
\newcommand{\phigh}{\mathbf{h}}
\newcommand{\epsaff}{\varepsilon_{\mathrm{aff}}}
\newcommand{\vansub}{\mathcal{V}}
\newcommand{\Wn}[1]{\widehat{\vansub}_{#1}}
\newcommand{\novel}{\mathcal{X}}
\newcommand{\mpoly}{\mathcal{P}}

% this paper
\newcommand{\ev}{\operatorname{ev}}            % the point evaluation functional
\newcommand{\cand}{\mathcal{T}}                % the set of candidates (challenge, list member of the fold)
\newcommand{\cls}{\mathcal{C}}                 % a class of candidates
\newcommand{\Exc}{\operatorname{Exc}}          % the exceptional set of point tuples
\newcommand{\pot}{\Psi}                        % the potential of a level
\newcommand{\thr}{\Theta}                      % the triple threshold K_0 (K_0 - 1)(K_0 - M)
\newcommand{\Kz}{K_0}                          % an integer constant of the analysis
\newcommand{\dsc}{\tau}                        % the discount of the potential
\newcommand{\epp}{\epsilon}                    % the error of one point
\newcommand{\Lb}{\mathrm{ListBound}}
\newcommand{\HK}{\mathrm{H}_{\K}}
\newcommand{\HE}{\mathrm{H}_{\E}}
\newcommand{\pull}{\operatorname{pull}}
\newcommand{\Bad}{\mathrm{Bad}}
\newcommand{\treeN}{\mathsf{N}}
\newcommand{\loc}{L}                           % a subspace polynomial L_V
\newcommand{\Zs}{Z}                            % a locator Z_S (context separates it from the oracle)
\newcommand{\Prob}{\Pr}
\newcommand{\Ex}{\mathbb{E}}
\newcommand{\bP}{\mathbf{P}}
\newcommand{\bV}{\mathbf{V}}
\DeclareMathOperator{\agrset}{agr}

% Lean names and status labels
\DeclareUrlCommand\leanurl{\urlstyle{tt}}
\newcommand{\lean}[1]{\leanurl{#1}}
\newcommand{\leanfile}[1]{\leanurl{lean/#1}}
\newcommand{\leansub}[2]{\leanurl{#1}\textsubscript{\texttt{#2}}}
\newcommand{\stmc}{\textsc{machine-checked}}
\newcommand{\stpr}{\textsc{proved}}
\newcommand{\stcj}{\textsc{conjectured}}
\newcommand{\status}[1]{\par\smallskip\noindent{\small\sloppy\textbf{Status.} #1\par}\smallskip}

% theorem environments
\theoremstyle{plain}
\newtheorem{theorem}{Theorem}[section]
\newtheorem{lemma}[theorem]{Lemma}
\newtheorem{corollary}[theorem]{Corollary}
\newtheorem{proposition}[theorem]{Proposition}
\newtheorem{conjecture}[theorem]{Conjecture}
\newtheorem{hypothesis}[theorem]{Hypothesis}
\theoremstyle{definition}
\newtheorem{definition}[theorem]{Definition}
\newtheorem{protocol}[theorem]{Protocol}
\newtheorem{problem}[theorem]{Open problem}
\theoremstyle{remark}
\newtheorem*{remark}{Remark}

% Lean listings
\lstdefinelanguage{lean}{
  morekeywords={theorem,lemma,def,noncomputable,by,fun,if,then,else,let,have,show,from,open,variable,abbrev},
  sensitive=true,
  morecomment=[l]{--},
}
\lstset{
  language=lean,
  basicstyle=\ttfamily\footnotesize,
  keywordstyle=\bfseries,
  columns=fullflexible,
  keepspaces=true,
  breaklines=true,
  breakindent=2em,
  frame=single,
  framesep=4pt,
  xleftmargin=6pt,
  xrightmargin=6pt,
  extendedchars=true,
  literate=
    {ℕ}{{$\mathbb{N}$}}1 {ℚ}{{$\mathbb{Q}$}}1 {∞}{{$\infty$}}1
    {→}{{$\to$}}1 {↔}{{$\leftrightarrow$}}1 {∀}{{$\forall$}}1 {∃}{{$\exists$}}1
    {≤}{{$\le$}}1 {≥}{{$\ge$}}1 {≠}{{$\ne$}}1 {¬}{{$\neg$}}1 {∧}{{$\wedge$}}1 {∨}{{$\vee$}}1
    {∈}{{$\in$}}1 {∉}{{$\notin$}}1 {⊆}{{$\subseteq$}}1 {∑}{{$\sum$}}1 {∏}{{$\prod$}}1
    {⊕}{{$\oplus$}}1 {×}{{$\times$}}1 {ˢ}{{$^{s}$}}1 {•}{{$\bullet$}}1 {∘}{{$\circ$}}1
    {⁻}{{$^{-}$}}1 {¹}{{$^{1}$}}1
    {₀}{{$_0$}}1 {₁}{{$_1$}}1 {₂}{{$_2$}}1 {₃}{{$_3$}}1 {ₗ}{{$_{l}$}}1
    {α}{{$\alpha$}}1 {β}{{$\beta$}}1 {σ}{{$\sigma$}}1 {τ}{{$\tau$}}1 {ι}{{$\iota$}}1 {κ}{{$\kappa$}}1
    {ζ}{{$\zeta$}}1 {φ}{{$\varphi$}}1 {Λ}{{$\Lambda$}}1 {ℓ}{{$\ell$}}1 {λ}{{$\lambda$}}1 {θ}{{$\theta$}}1
    {Γ}{{$\Gamma$}}1 {𝒞}{{$\mathcal{C}$}}1 {𝒯}{{$\mathcal{T}$}}1 {ε}{{$\varepsilon$}}1 {π}{{$\pi$}}1 {Θ}{{$\Theta$}}1
    {⟨}{{$\langle$}}1 {⟩}{{$\rangle$}}1 {↑}{{$\uparrow$}}1 {⊢}{{$\vdash$}}1 {∣}{{$\mid$}}1
}
